\documentclass[11pt,a4paper]{ctexart}

\usepackage[margin=1in]{geometry}
\usepackage{setspace}
\usepackage{booktabs}
\usepackage{longtable}
\usepackage{array}
\usepackage{graphicx}
\usepackage{amsmath,amssymb}
\usepackage{enumitem}
\usepackage{hyperref}
\usepackage{natbib}
\usepackage{titlesec}

\setstretch{1.2}
\hypersetup{
  colorlinks=true,
  linkcolor=black,
  citecolor=blue,
  urlcolor=blue
}

\titleformat{\section}{\large\bfseries}{\thesection}{0.75em}{}
\titleformat{\subsection}{\normalsize\bfseries}{\thesubsection}{0.75em}{}

\title{学习 AI 的角色：人类如何形成稳定的人机协作架构\\
\large Learning AI's Role: How Humans Converge Toward Stable Human-AI Collaboration Architectures}
\author{研究计划学术写作包}
\date{\today}

\begin{document}

\maketitle

\begin{abstract}
生成式人工智能进入决策系统后，组织真正需要回答的问题，不再只是是否使用 AI，而是如何配置 AI 在系统中的角色。现有研究虽然已经从 trust in automation、algorithm aversion、appropriate reliance 和 confidence calibration 等角度解释了人类如何看待和使用 AI，但它们大多聚焦于单次判断、局部信号或短期反应，尚不足以解释一个更动态的问题：人类如何在持续反馈中学习 AI 的系统角色，并最终收敛到某种稳定的人机协作架构。本文提出一个从 trust/adoption 转向 architecture learning 的研究框架，主张 AI error 的关键后果不是短期 trust 的下降，而是通过非对称的 attraction updating 改变长期收敛路径。为检验这一机制，本文提出一个基于多轮行为实验与 EWA-inspired 动态模型的研究设计，通过操控 AI error 与 human error 的时间分布，识别 early AI error 是否会导致低 AI autonomy 的路径依赖式锁定。本文预期在理论上推进 human--AI collaboration 研究的分析层级，在方法上将 EWA 从附加统计工具转化为全文的解释语言，并在实践上为 AI deployment 和 governance 提供关于 early-stage error management 的启发。
\end{abstract}

\tableofcontents
\newpage

\section{Introduction}

Generative AI is reshaping decision systems by changing not only how information is produced, but also how responsibility, oversight, and authority are distributed between humans and machines. In many contemporary domains, the critical organizational question is no longer whether AI should be used at all, but what role it should occupy in the decision pipeline. An AI system may provide advice, recommend a decision for human confirmation, act as a default that humans can override, or operate with substantial autonomy. These are not merely interface choices. They are institutional choices about delegation, supervision, accountability, and governance.

Existing scholarship has generated important insights into human responses to AI and automation. Research on trust in automation emphasizes that the key challenge is not blind acceptance but appropriate reliance \citep{lee2004trust}. Work on algorithm aversion shows that people may penalize algorithms heavily after witnessing errors \citep{dietvorst2015algorithm}, whereas work on algorithm appreciation suggests that under some task conditions people prefer algorithmic over human judgment \citep{logg2019algorithm}. A related body of research on AI-assisted decision making investigates how confidence scores, explanations, and performance cues calibrate trust and improve local decision quality \citep{zhang2020confidence,ma2023trust}. Yet these literatures largely analyze pointwise reactions: whether a system is trusted, whether advice is adopted, or whether case-level reliance is calibrated.

This paper argues that the more consequential problem lies one level higher. Humans do not only learn whether AI is reliable. They also learn what role AI should occupy in a decision system. The outcome of this learning process is not simply a revised belief about accuracy, but a revised collaboration architecture. Over repeated feedback, people may converge toward stable patterns of low or high AI authorization. Crucially, these trajectories may be path dependent: early AI errors may trigger stronger negative updating than comparable human errors, reducing later willingness to expose the system to AI-led configurations and thereby slowing recovery even when long-run AI performance equals that of human decision makers.

The central claim of this research prospectus is therefore straightforward: the key consequence of AI errors is not merely lower short-run trust, but altered convergence paths in human--AI collaboration architecture learning. To formalize this argument, the paper adapts the logic of experience-weighted attraction learning (EWA) \citep{camerer1999ewa}. Here, the strategic objects are not game-theoretic moves in the conventional sense, but collaboration architectures that differ in autonomy, supervision, cost, and risk exposure. This move allows us to connect repeated feedback, asymmetric learning from AI versus human errors, path dependence, architecture convergence, and behavioral steady states within a single dynamic framework.

The project contributes to ongoing debates in human--AI collaboration by shifting the analytical lens from adoption to role formation, from single advice episodes to long-run architecture convergence, and from learning about AI capability to learning about AI role. The empirical design proposed below is built to detect precisely these dynamics.

\section{Literature Review and Research Gap}

\subsection{From Trust to Architecture}

The foundational trust-in-automation literature established that performance in sociotechnical systems depends on appropriate reliance rather than simple acceptance or rejection \citep{lee2004trust}. This insight remains indispensable, but its typical unit of analysis is a particular automation interaction or local reliance decision. In AI-enabled organizations, however, actors also decide where AI should sit in the workflow, how much authority it should receive, and what forms of override or confirmation should remain in place. These are architecture questions, not merely trust questions.

\subsection{Algorithm Aversion, Appreciation, and the Limits of One-shot Framing}

Research on algorithm aversion demonstrates that people may avoid algorithms after observing them err, even when the algorithm performs better overall \citep{dietvorst2015algorithm}. By contrast, algorithm appreciation research shows that people may prefer algorithmic judgment under some circumstances \citep{logg2019algorithm}. Together, these findings imply that humans do not evaluate AI output in a source-neutral way. Still, the dominant framing remains one-shot: whether an algorithm is accepted, weighted, or rejected at a given moment. The present project asks a stronger question: can repeated asymmetries in how AI errors are learned produce long-run convergence toward low-authorization architectures?

\subsection{Confidence Calibration Is Not Role Learning}

Another important stream studies how confidence scores, explanations, and uncertainty signals affect trust calibration in AI-assisted decision making \citep{zhang2020confidence,ma2023trust}. This literature is chiefly concerned with signal-level learning: whether users learn the mapping between case-specific AI signals and correctness likelihood. That question is analytically distinct from the one pursued here. Even perfectly calibrated confidence signals do not by themselves specify whether AI should act autonomously, advise a human, or require explicit confirmation. The present research therefore differentiates signal-level calibration from architecture-level learning.

\subsection{Human--AI Teaming and Learning to Defer}

Recent research in machine learning and HCI has begun to consider how human and AI agents should divide labor. Learning-to-defer approaches ask when models should pass a decision to human experts \citep{mozannar2020defer}, and studies of human--AI teams examine how complementarity, performance tradeoffs, and updating affect joint outcomes \citep{bansal2019updates,vaccaro2024combinations}. These contributions are highly relevant, but they often approach the allocation problem from the system or model side. By contrast, this project foregrounds the human side of architecture formation: how people infer an appropriate role for AI under repeated feedback.

\subsection{EWA and Path Dependence as Integrating Mechanisms}

EWA provides a natural dynamic language for the proposed argument because it explains how repeated experience updates the attraction of alternative strategies \citep{camerer1999ewa}. Path dependence research, meanwhile, highlights how early events may shape long-run outcomes through lock-in and increasing returns \citep{david1985qwerty,arthur1989lockin}. Combining these traditions suggests a new mechanism for human--AI collaboration: early AI errors can induce stronger negative updating, leading actors to choose lower-autonomy structures, which in turn reduce future exposure to potentially trust-repairing AI performance, thereby slowing recovery and stabilizing low-AI-authorization equilibria in the loose behavioral sense.

\subsection{The Core Gap}

The literature therefore leaves three linked questions underdeveloped. First, how do humans learn AI's role, not just AI's reliability? Second, can asymmetric reactions to AI versus human errors shape long-run authorization structures rather than immediate preferences alone? Third, what formal mechanism can connect repeated feedback, architecture attraction, path dependence, and behavioral steady states? The proposed project is designed to answer these questions.

\section{Theory and Hypotheses}

\subsection{Core Theoretical Claim}

The project starts from a simple but consequential proposition: humans learn not only whether AI is good or bad, but where AI belongs in a decision system. This means that the relevant object of learning is the collaboration architecture itself. We define a human--AI collaboration architecture as the role allocation structure that specifies who decides, who recommends, who confirms, and who can override.

\subsection{Architecture Attraction}

Each architecture carries an ``attraction'' value: a subjective expectation of its desirability given prior feedback about accuracy, error source, correction opportunities, time cost, review burden, and perceived risk. Repeated experience updates these attractions. When the attractions of some architectures rise relative to others, later choices increasingly concentrate on them. Convergence occurs when this choice distribution stabilizes.

\subsection{Asymmetric Updating}

The key mechanism is asymmetric updating. If AI errors and human errors are learned symmetrically, then under equal long-run performance, early sequence differences should wash out over time. But if AI errors trigger stronger negative updates than comparable human errors, then early AI failures may disproportionately suppress later willingness to choose AI-led architectures. This introduces a self-reinforcing dynamic: lower authorization creates less later exposure to potentially trust-repairing AI performance, which slows recovery and increases the chance of long-run lock-in.

\subsection{Path Dependence and Behavioral Steady States}

The project does not claim a Nash equilibrium in the strict game-theoretic sense because the AI system is not modeled as an independently strategic actor. Instead, the paper uses the safer and more behaviorally appropriate language of behavioral steady states, learning equilibrium, and architecture convergence. The central concern is whether a repeated learning process stabilizes in a suboptimal authorization regime.

\subsection{Hypotheses}

\textbf{H1: Asymmetric negative learning.} When AI and human decision makers exhibit equivalent objective long-run performance, AI errors will trigger stronger negative updating of collaboration architectures than comparable human errors.

\textbf{H2: Path-dependent convergence.} Early AI errors will increase the likelihood that participants converge toward lower-AI-autonomy architectures.

\textbf{H3: Recovery asymmetry.} The positive effect of later correct AI performance on restoring AI authorization will be weaker than the negative effect of early AI errors on reducing AI authorization.

\textbf{H4: Human-error tolerance.} Early human errors will not symmetrically push participants toward substantially higher-AI-autonomy architectures.

\textbf{H5: Suboptimal steady states.} Even when AI and human long-run objective performance is equivalent, participants may settle into collaboration architectures that are more supervision-heavy than the objectively most efficient configuration.

\section{Research Design}

\subsection{Overview}

To test the argument, the project proposes a repeated-choice behavioral experiment in which participants act as supervisors of a decision pipeline. In each round, participants select one of five collaboration architectures:

\begin{enumerate}[label=\arabic*.]
    \item Human-only
    \item AI advice only
    \item AI recommendation + human confirmation
    \item AI default + human override
    \item AI autonomous decision
\end{enumerate}

These architectures can also be represented as an ordinal AI autonomy scale from 1 to 5.

\subsection{Core Identification Principle}

The central identification requirement is that AI and human analysts must have equal long-run objective performance. For example, both can be calibrated to an 80\% accuracy rate across the full experiment. The experimental manipulation then shifts only the timing and distribution of errors, not the long-run average quality of the human or AI decision source.

\subsection{Experimental Conditions}

Four conditions are proposed:

\begin{enumerate}[label=\arabic*.]
    \item \textbf{AI Early Error}: AI errors are concentrated in early rounds.
    \item \textbf{Human Early Error}: human errors are concentrated in early rounds.
    \item \textbf{Evenly Distributed Errors}: AI and human errors are distributed evenly across rounds.
    \item \textbf{AI Late Error}: AI errors are concentrated in later rounds.
\end{enumerate}

This design isolates whether error source and error timing affect later architecture choice trajectories.

\subsection{Round Structure}

Each round follows the same logic:

\begin{enumerate}[label=\arabic*.]
    \item The participant chooses a collaboration architecture.
    \item The system executes the chosen workflow.
    \item The participant receives feedback on correctness, gains/losses, time cost, review cost, whether the error was intercepted, and whether the error originated from AI or the human analyst.
    \item The participant proceeds to the next round.
\end{enumerate}

The main study should ideally contain around 80 rounds, with pretests determining whether 60 or 100 rounds would yield better convergence visibility without unacceptable fatigue.

\subsection{Dependent Variables}

The main dependent variables are architecture-level rather than trust-level:

\begin{itemize}[leftmargin=1.5em]
    \item final converged architecture,
    \item average AI autonomy level in late rounds,
    \item probability of entering a low-AI-authorization steady state,
    \item convergence speed,
    \item recovery speed after AI-error shocks,
    \item choice volatility in late rounds,
    \item distance between the stable architecture and the objectively optimal architecture.
\end{itemize}

\subsection{Steady-State Criteria}

Behavioral steady states can be defined using multiple parallel criteria:

\begin{itemize}[leftmargin=1.5em]
    \item a single architecture accounts for more than 70\% of choices in the final 15 rounds,
    \item variance in autonomy level during the final 15 rounds falls below a predefined threshold,
    \item architecture choice probabilities no longer change significantly over time,
    \item estimated attractions from the structural model enter a stable band.
\end{itemize}

\subsection{EWA-inspired Model}

Let $A_{k,t}$ denote the attraction of architecture $k$ after round $t$, and let $s_t$ denote the chosen architecture in round $t$. The updating logic follows an EWA-style process:

\begin{align}
N_t &= \rho N_{t-1} + 1 \\
A_{k,t} &= \frac{\phi N_{t-1}A_{k,t-1} + [\delta + (1-\delta)I(s_t = k)]U_{k,t}}{N_t} \\
P(s_{t+1}=k) &= \frac{\exp(\beta A_{k,t})}{\sum_j \exp(\beta A_{j,t})}
\end{align}

To capture asymmetric learning, the instantaneous utility term can be decomposed by feedback source and sign:

\begin{equation}
U_{k,t} =
\alpha_{AI}^{+} I(AI\ correct)
+ \alpha_{AI}^{-} I(AI\ error)
+ \alpha_{H}^{+} I(Human\ correct)
+ \alpha_{H}^{-} I(Human\ error)
- c_t
\end{equation}

The critical parameter tests are:

\begin{align}
\alpha_{AI}^{-} &> \alpha_{H}^{-} \\
|\alpha_{AI}^{-}| &> |\alpha_{AI}^{+}|
\end{align}

\subsection{Analysis Strategy}

The empirical analysis proceeds in three layers. First, descriptive trajectory plots show how architecture use and autonomy change across rounds by condition. Second, multilevel or GEE models estimate the long-run effect of experimental conditions on later authorization. Third, the EWA-inspired model is estimated structurally to identify whether asymmetric learning weights and recovery asymmetries account for the observed path dependence.

\section{Expected Contributions}

\subsection{Theoretical Contribution}

The first contribution is conceptual. The project reframes the problem of human interaction with AI from trust or adoption to architecture learning. This shifts the unit of analysis from immediate reliance decisions to evolving role structures within a decision system.

\subsection{Mechanism Contribution}

The second contribution is mechanistic. Rather than claiming only that AI errors reduce trust, the project shows how AI errors may alter long-run convergence paths by generating stronger negative attraction updates than equivalent human errors. This turns algorithm aversion into a path-dependent learning problem rather than a momentary preference anomaly.

\subsection{Methodological Contribution}

The third contribution is methodological. By integrating repeated behavioral experiments with an EWA-inspired structural model, the project links experimental manipulation, behavioral trajectories, and interpretable learning parameters. In this framework, EWA is not an afterthought; it is the theoretical language of the project.

\subsection{Practical Contribution}

The fourth contribution concerns AI governance. If early AI errors shape long-run authorization regimes, then deployment strategy matters as much as average model quality. Organizations may need staged delegation, reversible autonomy, buffered pilot phases, and explicit error-repair protocols to prevent low-automation lock-in driven by early shocks.

\section{Conclusion}

This project argues that the key research question in the AI era is no longer simply whether humans trust AI, adopt AI, or accept a particular AI recommendation. The deeper question is how humans learn what role AI should occupy in a decision system and how repeated feedback shapes the long-run architecture of collaboration. By combining architecture-level theorizing, path-dependence logic, and EWA-inspired formalization, the proposed research offers a more dynamic account of human--AI collaboration than existing one-shot trust or reliance frameworks. Its ultimate value lies in showing that the consequences of AI error may be organizationally persistent, not merely psychologically immediate.


\bibliographystyle{plainnat}
\bibliography{references}

\end{document}
